\section{总结与延伸}

\subsection{核心结论总表}

\begin{table}[H]
\centering
\caption{访谈核心观点与可执行启示}
\begin{tabular}{p{0.22\textwidth} p{0.34\textwidth} p{0.34\textwidth}}
\toprule
主题 & Jensen 的关键观点 & 对工程团队的启示 \\
\midrule
系统边界 & 竞争从 chip scale 走向 rack scale 与 AI factory & KPI 要从单机性能升级为系统吞吐、稳定性和交付效率 \\
\addlinespace
生态护城河 & CUDA 的 installed base 是首要优势 & 平台团队应优先建设工具链、文档、社区和迁移路径 \\
\addlinespace
产品演进 & Grace Blackwell 到 Vera Rubin 反映 workload 驱动迭代 & 路线图需要绑定真实负载数据，而非只按发布节奏推进 \\
\addlinespace
供应链能力 & 200 家合作方与工业化制造决定交付上限 & 架构设计需前置考虑制造、测试、运输和运维约束 \\
\addlinespace
能源约束 & AI 扩张受 power grid 和能耗效率约束 & 规划阶段把电力与散热当成一等公民，建立能耗治理机制 \\
\addlinespace
领导力机制 & 组织结构应匹配跨层协同复杂度 & 采用跨域问题评审与快速复盘，减少串行信息损耗 \\
\bottomrule
\end{tabular}
\end{table}

\begin{importantbox}{一句话总结}
这场访谈最重要的启示是：AI 时代的长期优势不在某一代芯片，而在 ``系统工程能力 + 生态飞轮 + 产业协同'' 的复合体。任何只优化单点的策略，都会在规模化阶段暴露上限。
\end{importantbox}

\subsection{可延伸阅读}

\begin{itemize}
    \item Lex Fridman Podcast \#494（原视频）：\url{https://www.youtube.com/watch?v=vif8NQcjVf0}
    \item NVIDIA Data Center / NVLink 资料：\url{https://www.nvidia.com/en-us/data-center/nvlink/}
    \item CUDA Documentation：\url{https://docs.nvidia.com/cuda/}
    \item NVIDIA Grace Blackwell 平台介绍：\url{https://www.nvidia.com/en-us/data-center/grace-blackwell-superchip/}
\end{itemize}

%% --- Fin del contenido --- %%

\end{document}
